\section{Scope, antecedents, and the main mechanism}\label{shj:scope}
The canonical ProveIt manuscript connects polylogarithms, gamma functions, Hurwitz-zeta jets, and Stieltjes antiderivatives. Its integration and differentiation chapters already contain the pointwise ladders relating these objects. The present continuation introduces a bilinear operation: translate one function around the unit circle and integrate its product with another. Translation is important. For $0<a<1$, the two endpoint singularities occur at different points, allowing a canonical finite-part product without multiplying distributions at a common singularity.

The governing observation is that the shifted Hurwitz correlation has two exact representations. Fourier multiplication produces polylogarithms of order $2-s-t$. A functional equation then replaces those polylogarithms by Hurwitz zeta of order $s+t-1$. Near $(s,t)=(1,1)$, a gamma quotient divisible by $u(u+v)$ controls all finite coefficients. This divisibility explains both the linear Stieltjes closure and the absence of one apparently possible Stieltjes index.

The Fourier expansion, gamma identities, and basic Hurwitz calculus are classical; standard conventions are recorded in the NIST Digital Library of Mathematical Functions \cite{shj:DLMF25,shj:DLMFpoly}. Espinosa and Moll treat ordinary same-point and reflected Hurwitz products and their differentiated integrals \cite{shj:EM}. Coffey develops Stieltjes and Hurwitz-zeta representations \cite{shj:Coffey}. Spectral operator interpretations of the Hurwitz and polylogarithm families also appear in Gay and Sebbar \cite{shj:GS}. We give the shifted two-variable identity a self-contained proof rather than presenting the Fourier method as a new discovery.

The specific contribution of this report is the explicit, endpoint-normalized package: the all-index Stieltjes coefficient formula; its harmonic contact correction under differentiation; and its systematic lift to convergent antiderivative correlations. These statements are proved below, not inferred from integer-relation searches. A literature search located the classical antecedents just mentioned but did not establish global priority for every specialization. Accordingly, ``continuation'' means a developed addition to the inspected repository material, not an assertion that no equivalent formula exists elsewhere.

The inspected sources and observed blob identifiers are recorded in \src{data/source_snapshot.json}. The current manuscript explicitly keeps $S_6$ and $S_8$ conjectural and distinguishes small residuals from proofs \cite{shj:repo}. Those problems are not solved here. We also do not infer numerical independence of Hurwitz or Dirichlet coordinates. The reflected \emph{powers} of $\log\Gamma$ studied elsewhere in the manuscript are different from the shifted, bilinear spectral-derivative correlations below.

\section{Conventions and endpoint normalization}\label{shj:conventions}
Throughout,
\[
0<a<1,\qquad b=1-a,\qquad \ell=\log(2\pi),\qquad
\kappa=\gamma+\ell,\qquad B_2(a)=a^2-a+\frac16.
\]
The fractional part $\{x+a\}$ is interpreted periodically; its value at the single jump point does not affect ordinary integrals. The notation
\[
\zeta^{[j]}(s,x)=\partial_s^j\zeta(s,x),\qquad
\Li_s^{[j]}(z)=\partial_s^j\Li_s(z)
\]
always means \emph{spectral} differentiation. In contrast, $\gamma_m^{(r)}(x)$ and $\psi^{(r)}(x)$ mean derivatives in the real argument. Euler's constant is $\gamma=\gamma_0(1)$, not a new parameter.

The generalized Stieltjes constants are normalized by
\begin{equation}\label{shj:laurent}
\zeta(1+u,x)=\frac1u+R(u,x),\qquad
R(u,x)=\sum_{m\ge0}\frac{(-1)^m}{m!}\gamma_m(x)u^m.
\end{equation}
Thus $\gamma_0(x)=-\psi(x)$. The Hurwitz shift identity gives the exact local decomposition
\begin{equation}\label{shj:localgamma}
\gamma_m(x)=\frac{\log^m x}{x}+\gamma_m(1+x),\qquad x>0.
\end{equation}
Every logarithm of a positive real argument is the real logarithm. For $z=e^{2\pi\ii a}\ne1$, $\Li_s(z)$ and $\Li_s(z^{-1})$ denote their usual boundary values from the unit disk, continued in $s$; these are entire in $s$. For real $s$ they are conjugate, but we retain them separately when $s$ is complex.

\begin{definition}[Canonical shifted Stieltjes finite part]\label{shj:defJ}
For nonnegative integers $m,n$, define
\begin{align}
J_{mn}(a)=\lim_{\epsilon\downarrow0}\bigg\{&
\int_\epsilon^b\gamma_m(x)\gamma_n(x+a)\dd x
+\int_\epsilon^a\gamma_m(b+t)\gamma_n(t)\dd t\notag\\
&+\frac{\gamma_n(a)}{m+1}\log^{m+1}\epsilon
+\frac{\gamma_m(b)}{n+1}\log^{n+1}\epsilon\bigg\}.
\label{shj:Jcutoff}
\end{align}
Equivalently, this is denoted
$\FP\int_0^1\gamma_m(x)\gamma_n(\{x+a\})\dd x$.
\end{definition}
The two cutoffs can tend to zero independently: each counterterm belongs to its own local endpoint. The coordinates are exactly $x$ and $t=x-b$, with no rescaling or additional subtraction length. A different finite-part convention can change local constants and therefore change the identities.

\begin{lemma}[Existence and an ordinary-integral version]\label{shj:Jexists}
The limit \eqref{shj:Jcutoff} exists and is equivalent to
\begin{align}
J_{mn}(a)={}&\int_0^b\left[
\gamma_m(x)\gamma_n(x+a)-\gamma_n(a)\frac{\log^m x}{x}\right]\dd x\notag\\
&+\int_0^a\left[
\gamma_m(b+t)\gamma_n(t)-\gamma_m(b)\frac{\log^n t}{t}\right]\dd t\notag\\
&+\frac{\gamma_n(a)}{m+1}\log^{m+1}b
+\frac{\gamma_m(b)}{n+1}\log^{n+1}a.
\label{shj:Jregular}
\end{align}
Both integrals on the right are ordinary improper integrals. Moreover,
$J_{mn}(a)=J_{nm}(1-a)$.
\end{lemma}
\begin{proof}
By \eqref{shj:localgamma}, subtracting the singular term in the first integral leaves $O(1+|\log x|^m)$ near zero: the other factor differs from $\gamma_n(a)$ by $O(x)$. The second integral has the same property. Integrating $\log^m x/x$ explicitly yields \eqref{shj:Jregular}. Interchanging the two pieces and replacing $a$ by $b$ proves the symmetry.
\end{proof}
We use analogous finite parts for polygamma and derivative products: expand each smooth factor at the distinct singular point, integrate the resulting powers and logarithms, and retain the constant term after deleting all divergent cutoff terms. Section~\ref{shj:contacts} makes this prescription precise through periodic distributions.

\section{The shifted Hurwitz correlation}\label{shj:master}
\begin{proposition}[Classical Fourier method, shifted form]\label{shj:masterthm}
For $\Re s<1$, $\Re t<1$, define the ordinary integral
\begin{equation}\label{shj:Cdef}
C(s,t;a)=\int_0^1\zeta(s,x)\zeta(t,\{x+a\})\dd x.
\end{equation}
It is jointly holomorphic in this domain. Put $w=2-s-t$ and $z=e^{2\pi\ii a}$. Then
\begin{align}
C(s,t;a)=\frac{\Gamma(1-s)\Gamma(1-t)}{(2\pi)^w}
\bigg[&e^{-\pi\ii(s-t)/2}\Li_w(z)\notag\\
&+e^{\pi\ii(s-t)/2}\Li_w(z^{-1})\bigg].
\label{shj:Cpoly}
\end{align}
It also has the meromorphic representation
\begin{align}
C(s,t;a)={}&B(1-s,s+t-1)\zeta(s+t-1,a)\notag\\
&+B(1-t,s+t-1)\zeta(s+t-1,b).
\label{shj:Cbeta}
\end{align}
At exceptional arguments of the beta factors, the \emph{sum} in \eqref{shj:Cbeta} is continued before evaluation. Apparent poles of separate summands need not be poles of $C$.
\end{proposition}
\begin{proof}
Split the integral at $b$. At $x=0$ only the first factor is singular and has the form $x^{-s}$ plus a smooth function. At $x=b+$ only the second factor is singular. Compact subsets of $\Re s<1$, $\Re t<1$ therefore admit integrable dominating functions, including arbitrary spectral derivatives, because powers of $|\log x|$ remain integrable against $x^{-\sigma}$ for $\sigma<1$. This proves convergence and holomorphy. There is no additional condition $\Re(s+t)<1$ for a nonzero shift.

For $\Re s<0$, the Hurwitz Fourier expansion \cite[Eq.~25.11.9]{shj:DLMF25} is absolutely convergent and its nonzero coefficients are
\begin{equation}\label{shj:fouriercoef}
\widehat\zeta_s(k)=\Gamma(1-s)(2\pi|k|)^{s-1}
\exp\!\left(-\frac{\pi\ii}{2}(1-s)\sgn k\right),\qquad k\ne0.
\end{equation}
The zero coefficient vanishes. For $\Re s,\Re t<0$, multiply the two absolutely convergent series and integrate. Orthogonality gives
\[
C(s,t;a)=\sum_{k\ne0}\widehat\zeta_s(k)\widehat\zeta_t(-k)e^{-2\pi\ii ka}.
\]
The positive and negative frequencies give exactly \eqref{shj:Cpoly}. Both sides are holomorphic on the larger connected domain stated in the proposition; the identity theorem extends the equality there.

To obtain \eqref{shj:Cbeta}, apply the same Fourier formula to $\zeta(1-w,a)$ and $\zeta(1-w,b)$, initially away from gamma poles. Euler reflection gives
$\Gamma(w)\Gamma(1-w)=\pi/\sin\pi w$ and
$1/\Gamma(t)=\Gamma(1-t)\sin\pi t/\pi$.
The coefficient of $\Li_w(z)$ reduces using
\[
\frac{\sin(\pi t)e^{-\pi\ii w/2}+\sin(\pi s)e^{\pi\ii w/2}}
{\sin(\pi w)}=e^{-\pi\ii(s-t)/2},\qquad w=2-s-t.
\]
The other coefficient is obtained by reversing the phase. This proves equality on an open set and hence meromorphically everywhere.
\end{proof}

\begin{corollary}[All ordinary spectral-jet correlations]\label{shj:ordinaryjets}
For integers $p,q\ge0$ and $\Re s,\Re t<1$,
\[
\int_0^1\zeta^{[p]}(s,x)\zeta^{[q]}(t,\{x+a\})\dd x
=\partial_s^p\partial_t^q C(s,t;a).
\]
Thus \eqref{shj:Cpoly} generates ordinary identities at nonpositive spectral integers, while continuation of \eqref{shj:Cbeta} generates finite-part identities at positive integers.
\end{corollary}
\begin{proof}
Differentiate under the integral using the compact domination established above. The formulas on the right are finite applications of the product and chain rules.
\end{proof}
The distinction between nonzero and zero shifts is essential. At $a=0$, the singularities coincide and the original convergence region changes to include a constraint on $s+t$. One must not substitute $a=0$ directly into a separated-singularity finite-part formula. The continuous ordinary log-gamma correlation will have a legitimate limit at zero, but that is a different assertion.

\section{Periodic continuation and the missing point mass}\label{shj:periodic}
Let $\T=\R/\Z$. A distribution on $\T$ acts on smooth periodic test functions $\phi$. Write $\delta_0$ for evaluation at zero and $D$ for distributional differentiation. For $\Re s<1$, let $Z_s$ be the periodic distribution represented on $(0,1)$ by $\zeta(s,x)$.

\begin{lemma}[Meromorphic periodic extension]\label{shj:periodicresidue}
The family $Z_s$ extends meromorphically to all $s\in\C$. Its residue at $s=1$ is $1-\delta_0$. At $s=r+1$ for $r\ge1$, its residue is
\begin{equation}\label{shj:residuepositive}
\Res_{s=r+1}Z_s=\frac{(-1)^{r+1}}{r!}\delta_0^{(r)}.
\end{equation}
For every $m\ge0$, define
\begin{equation}\label{shj:Gdef}
\langle G_m,\phi\rangle=
\lim_{\epsilon\downarrow0}\left[
\int_\epsilon^1\gamma_m(x)\phi(x)\dd x
+\frac{\phi(0)}{m+1}\log^{m+1}\epsilon\right].
\end{equation}
Then $G_m$ has mean zero and
\begin{equation}\label{shj:Zlaurent}
Z_{1+u}=\frac{1-\delta_0}{u}
+\sum_{m\ge0}\frac{(-1)^m}{m!}G_m u^m.
\end{equation}
\end{lemma}
\begin{proof}
Use $\zeta(s,x)=x^{-s}+\zeta(s,1+x)$. For any fixed $M$, subtract the first $M$ terms of the Taylor expansion of $\phi$ at zero inside the first integral. The subtracted terms integrate to
\[
\sum_{j=0}^{M-1}\frac{\phi^{(j)}(0)}{j!(j+1-s)},
\]
and the remainder is holomorphic for $\Re s<M+1$. Increasing $M$ gives the continuation. The residue of the $j=r$ term at $s=r+1$ is $-\phi^{(r)}(0)/r!$, which equals the action of the distribution in \eqref{shj:residuepositive}. The smooth term $\zeta(s,1+x)$ contributes an additional residue $1$ only at $s=1$.

Expanding $x^{-1-u}$ in the same subtraction formula proves \eqref{shj:Gdef} and \eqref{shj:Zlaurent}, using \eqref{shj:localgamma}. The zero Fourier coefficient of $Z_s$ vanishes for $\Re s<1$: integrate $\partial_x\zeta(s-1,x)=(1-s)\zeta(s,x)$ and use the shift relation at zero. It therefore vanishes meromorphically. The pole coefficient $1-\delta_0$ also has zero mean. Comparing regular coefficients gives $\langle G_m,1\rangle=0$.
\end{proof}
The familiar pointwise residue of $\zeta(s,x)$ is $1$; the periodic distributional residue is $1-\delta_0$. There is no contradiction: the second residue records the nonintegrable endpoint that appears when the pointwise family is integrated against a test function. Losing this point mass leads to incorrect differentiated finite-part formulas.

\begin{proposition}[Coefficient extraction at the double pole]\label{shj:Jcoefficient}
For $a\in(0,1)$,
\begin{equation}\label{shj:JfromC}
J_{mn}(a)=(-1)^{m+n}m!n![u^m v^n]C(1+u,1+v;a).
\end{equation}
The brackets mean the regular Laurent coefficient at the specified bidegree. In particular, the double-pole coefficient is $-1/(uv)$, not $1/(uv)$.
\end{proposition}
\begin{proof}
The singular supports of $Z_{1+u}(x)$ and $Z_{1+v}(x+a)$ are $0$ and $b$, which are distinct. Choose disjoint neighborhoods of these points and a smooth partition of unity. On each neighborhood only one factor is a distribution and the other is smooth, so their product and integral are defined without any product of distributions at a common singularity. This gives the meromorphic continuation of the original integral.

The product of the regular coefficients in \eqref{shj:Zlaurent}, integrated with this partition, is exactly \eqref{shj:Jcutoff}. The pole terms give
\begin{align}
C(1+u,1+v;a)={}&-\frac1{uv}-\frac{R(v,a)}u-\frac{R(u,b)}v\notag\\
&+\sum_{m,n\ge0}\frac{(-1)^{m+n}}{m!n!}J_{mn}(a)u^mv^n.
\label{shj:fullLaurent}
\end{align}
The integrated product of $1-\delta_0$ and $1-\delta_b$ is $1-1-1=-1$ away from $a=0$, and each $G_m$ has mean zero. This proves the statement.
\end{proof}

\begin{corollary}[Polylogarithmic generator]\label{shj:polygen}
Set $E(u)=\Gamma(1-u)(2\pi)^u$ and
\begin{align}
\mathcal H(u,v;a)=E(u)E(v)\big[&e^{-\pi\ii(u-v)/2}\Li_{-u-v}(z)\notag\\
&+e^{\pi\ii(u-v)/2}\Li_{-u-v}(z^{-1})\big].
\label{shj:H}
\end{align}
Then
\begin{equation}\label{shj:Jpoly}
J_{mn}(a)=(-1)^{m+n}m!n![u^{m+1}v^{n+1}]\mathcal H(u,v;a).
\end{equation}
\end{corollary}
\begin{proof}
In \eqref{shj:Cpoly}, set $s=1+u$, $t=1+v$ and use $\Gamma(-u)=-\Gamma(1-u)/u$. Thus $C=\mathcal H/(uv)$, and \eqref{shj:JfromC} applies.
\end{proof}
This representation explains how polylogarithm \emph{order derivatives}, rather than only argument derivatives, enter the exact Stieltjes product calculus.

\section{An all-order linear Stieltjes closure theorem}\label{shj:closure}
The beta representation is particularly effective near $(s,t)=(1,1)$. Define
\begin{equation}\label{shj:AandT}
A(u,v)=\frac{\Gamma(1-u)\Gamma(1+u+v)}{\Gamma(1+v)},
\qquad T(u,v)=\frac{A(u,v)-1}{u(u+v)}.
\end{equation}
The apparent singularities in $T$ are removable: $A(0,v)=1$ and $A(u,-u)=1$. Its Taylor coefficients belong to the polynomial algebra generated by positive integer zeta values. More explicitly,
\begin{equation}\label{shj:logA}
\log A(u,v)=\sum_{k\ge2}\frac{\zeta(k)}k
\left[u^k+(-1)^k\big((u+v)^k-v^k\big)\right].
\end{equation}
Euler's constant cancels identically. This series converges in a neighborhood of the origin; it can also be used as an exact formal power series.

\begin{theorem}[All-index finite-part closure]\label{shj:closurethm}
Let $m,n\ge0$ and $d=m+n$. Define the analytic germ
\begin{align}
\mathcal F(u,v;a)={}&T(u,v)+T(v,u)\notag\\
&+(u+v)\big[T(u,v)R(u+v,a)+T(v,u)R(u+v,b)\big].
\label{shj:Fclosure}
\end{align}
Then
\begin{equation}\label{shj:closureformula}
\boxed{\displaystyle
J_{mn}(a)=\frac{\gamma_{d+1}(a)}{m+1}
+\frac{\gamma_{d+1}(b)}{n+1}
-(-1)^d m!n![u^m v^n]\mathcal F(u,v;a).}
\end{equation}
Every coefficient is obtained by a finite calculation. Apart from the two leading terms displayed explicitly, the right side contains only $\gamma_j(a)$ and $\gamma_j(b)$ with $0\le j\le d-1$, and a constant term. Its coefficients lie in
$\Q[\zeta(2),\ldots,\zeta(d+2)]$.
There are no products of generalized Stieltjes functions and no term involving $\gamma_d$.
\end{theorem}
\begin{proof}
Since $B(-u,1+u+v)=-A(u,v)/u$, equation \eqref{shj:Cbeta} becomes
\begin{equation}\label{shj:Cnear}
C(1+u,1+v;a)=-\frac{A(u,v)}u\zeta(1+u+v,a)
-\frac{A(v,u)}v\zeta(1+u+v,b).
\end{equation}
Substitute $A(u,v)=1+u(u+v)T(u,v)$ and the Laurent expansion \eqref{shj:laurent}. The result is
\begin{equation}\label{shj:Csplit}
C(1+u,1+v;a)=-\frac1{uv}
-\frac{R(u+v,a)}u-\frac{R(u+v,b)}v-\mathcal F(u,v;a).
\end{equation}
Here the apparent diagonal pole at $u+v=0$ has canceled exactly.

The regular coefficient of $-R(u+v,a)/u$ at $u^mv^n$ is
\[
-\frac{(-1)^{d+1}\gamma_{d+1}(a)}{(d+1)!}
\binom{d+1}{m+1}
=\frac{(-1)^d\gamma_{d+1}(a)}{(m+1)!n!}.
\]
The other term similarly contributes
$(-1)^d\gamma_{d+1}(b)/(m!(n+1)!)$.
Now apply \eqref{shj:JfromC} to get \eqref{shj:closureformula}.

To obtain a coefficient of total degree $d$ in $T$, only terms of total degree at most $d+2$ in $A$ are required. Formula \eqref{shj:logA} proves the asserted zeta coefficient algebra. In the product $(u+v)TR$, a Stieltjes index $j$ occupies degree $j+1$ before multiplication by $T$, so $j\le d-1$. All appearances of $R$ are linear. These observations prove the remaining claims.
\end{proof}

\begin{corollary}[Formal weight conservation]\label{shj:weight}
Assign weight $j+1$ to $\gamma_j(a)$ and $\gamma_j(b)$ and weight $k$ to $\zeta(k)$. Every monomial in \eqref{shj:closureformula} has weight $m+n+2$.
\end{corollary}
\begin{proof}
The degree-$k$ component of $\log A$ has coefficient weight $k$. Thus a coefficient of total degree $r$ in $T$ has weight $r+2$. A contribution from $(u+v)TR$ to total degree $d$ has $r+j+1=d$ and weight $(r+2)+(j+1)=d+2$. The leading Stieltjes terms have the same weight.
\end{proof}
This is a grading of the displayed symbolic algebra, not a theorem that the evaluated constants have no additional relations.

\subsection{The first four nontrivial identities}
The beginning of the universal kernel is
\begin{equation}\label{shj:Tstart}
T(u,v)=\zeta(2)-\zeta(3)v
+\frac{\zeta(4)+\zeta(2)^2}{2}(u^2+uv)
+\zeta(4)v^2+O((|u|+|v|)^3).
\end{equation}
Write $S_j(a)=\gamma_j(a)+\gamma_j(b)$. Coefficient extraction gives
\begin{align}
J_{00}(a)&=S_1(a)-2\zeta(2),\label{shj:J00}\\
J_{10}(a)&=\tfrac12\gamma_2(a)+\gamma_2(b)
+\zeta(2)S_0(a)-\zeta(3),\label{shj:J10}\\
J_{11}(a)&=\tfrac12 S_3(a)+2\zeta(2)S_1(a)+\zeta(3)S_0(a)
-\zeta(4)-\zeta(2)^2,\label{shj:J11}\\
J_{20}(a)&=\tfrac13\gamma_3(a)+\gamma_3(b)+2\zeta(2)S_1(a)
+2\zeta(3)\gamma_0(b)\notag\\
&\hspace{2cm}-3\zeta(4)-\zeta(2)^2.\label{shj:J20}
\end{align}
These are exact analytic identities under Definition~\ref{shj:defJ}, not conjectured PSLQ relations. The asymmetric coefficients in $J_{10}$ and $J_{20}$ are necessary; exchanging $a$ and $b$ exchanges the two Laurent indices.

Because $\gamma_0=-\psi$, the first identity takes the particularly simple form
\begin{equation}\label{shj:psipsi}
\FP\int_0^1\psi(x)\psi(\{x+a\})\dd x
=\gamma_1(a)+\gamma_1(b)-\frac{\pi^2}{3}.
\end{equation}
An equivalent identity containing only ordinary integrals is
\begin{align}
&\int_0^b\left[\psi(x)\psi(x+a)+\frac{\psi(a)}x\right]\dd x
+\int_0^a\left[\psi(b+t)\psi(t)+\frac{\psi(b)}t\right]\dd t\notag\\
&\qquad=\gamma_1(a)+\gamma_1(b)-\frac{\pi^2}{3}
+\psi(a)\log b+\psi(b)\log a.
\label{shj:psipsireg}
\end{align}
The endpoint singularities cancel inside each bracket. This version is useful when a conventional improper integral is preferable to distribution notation.

\subsection{A finite exact coefficient algorithm}
For target total index $d$, truncate \eqref{shj:logA} at degree $d+2$. Exponentiate by
\[
A=\sum_{j=0}^{\lfloor(d+2)/2\rfloor}\frac{(\log A)^j}{j!}
\pmod{\text{total degree}>d+2}.
\]
Exact polynomial division of $A-1$ by $u(u+v)$ has zero remainder. Let
$T(u,v)=\sum t_{ij}u^iv^j$. The correction coefficient needed in \eqref{shj:closureformula} is
\begin{align}
[u^mv^n]\mathcal F={}&t_{mn}+t_{nm}\notag\\
&+\sum_{\substack{i\le m,\ j\le n\\k=d-i-j-1\ge0}}
 t_{ij}\frac{(-1)^k}{k!}\binom{k+1}{m-i}\gamma_k(a)\notag\\
&+\sum_{\substack{j\le m,\ i\le n\\k=d-i-j-1\ge0}}
 t_{ij}\frac{(-1)^k}{k!}\binom{k+1}{m-j}\gamma_k(b).
\label{shj:coefficientalgorithm}
\end{align}
Only polynomial arithmetic over rational numbers and formal zeta symbols is used. The implementation does not perform floating-point relation recognition. Exact divisibility, reflection symmetry, linearity, and weight conservation provide independently checkable invariants of the generated formulas.

\section{Differentiation and the contact-term theorem}\label{shj:contacts}
Let $T_{m,r}$ denote the canonical periodic finite-part extension of the \emph{pointwise} function $\gamma_m^{(r)}(x)$, with $T_{m,0}=G_m$. To define it on a test function, integrate from $\epsilon$ to $1$, expand the singular powers and logarithms using the Taylor coefficients of the test function at zero, remove all divergent terms, and retain the constant term. This is the same local-coordinate normalization used in Definition~\ref{shj:defJ}.

For $r,m\ge0$, put
\begin{equation}\label{shj:hdef}
h_{r,m}=(-1)^m m![u^{m+1}]\prod_{j=1}^r\left(1+\frac uj\right).
\end{equation}
The empty product is $1$. In particular,
\[
h_{0,m}=0,\qquad h_{r,0}=H_r,\qquad
h_{r,m}=(-1)^m m!e_{m+1}(1,1/2,\ldots,1/r),
\]
where $H_r=\sum_{j=1}^r1/j$ and $H_0=0$. Thus $h_{r,m}=0$ whenever $m\ge r$.

\begin{theorem}[Exact contact correction]\label{shj:contactthm}
As distributions on $\T$,
\begin{equation}\label{shj:contactlaw}
\boxed{D^rG_m=T_{m,r}-h_{r,m}\delta_0^{(r)}.}
\end{equation}
Consequently, for $r,k,m,n\ge0$ and $N=r+k$,
\begin{align}
&\FP\int_0^1\gamma_m^{(r)}(x)\gamma_n^{(k)}(\{x+a\})\dd x\notag\\
&\quad=(-1)^r J_{mn}^{(N)}(a)
+(-1)^k h_{k,n}\gamma_m^{(N)}(b)
+(-1)^r h_{r,m}\gamma_n^{(N)}(a).
\label{shj:generalcontact}
\end{align}
All terms on the right are ordinary functions for $0<a<1$.
\end{theorem}
\begin{proof}
For $\Re s$ sufficiently negative, integration by parts and the matching periodic endpoint values give
\[
D^rZ_s=(-1)^r(s)_rZ_{s+r},
\qquad (s)_r=s(s+1)\cdots(s+r-1).
\]
The identity holds meromorphically for every $s$. Set $s=1+u$ and write
$(1+u)_r=r!p_r(u)$, where $p_r(u)=\prod_{j=1}^r(1+u/j)$.
For $r\ge1$, the pole part on the right, from \eqref{shj:residuepositive}, is
\[
(-1)^r r!p_r(u)\frac{(-1)^{r+1}\delta_0^{(r)}}{r!u}
=-\frac{p_r(u)}u\delta_0^{(r)}.
\]
Its leading pole $-\delta_0^{(r)}/u$ matches $D^r(1-\delta_0)/u$ on the left. The remaining polynomial coefficients in this polar contribution are the contact terms.

More explicitly, the local singular factor is
$(-1)^r(1+u)_r x^{-r-1-u}$. After its pole is separated, coefficient extraction commutes with the prescribed constant-term extension of powers and logarithms. Thus the regular, nonresidue coefficient at $u^m$ is $(-1)^mT_{m,r}/m!$. The additional coefficient from $-(p_r(u)-1)\delta_0^{(r)}/u$ is $-[u^{m+1}]p_r(u)\delta_0^{(r)}$. Multiplying by $(-1)^m m!$ proves \eqref{shj:contactlaw}. The case $r=0$ is immediate.

For the correlation identity substitute
$T_{m,r}=D^rG_m+h_{r,m}\delta_0^{(r)}$ and the corresponding expression for $T_{n,k}$. The product of the two delta derivatives is zero for nonzero shift, because their supports are disjoint. Periodic integration by parts gives
\[
\int (D^rG_m)(x)(D^kG_n)(x+a)\dd x
=(-1)^r\frac{d^{r+k}}{da^{r+k}}J_{mn}(a).
\]
The delta at $b$ contributes $(-1)^k h_{k,n}\gamma_m^{(r+k)}(b)$, and the delta at zero contributes $(-1)^r h_{r,m}\gamma_n^{(r+k)}(a)$. These operations are legitimate using the disjoint-neighborhood partition from Proposition~\ref{shj:Jcoefficient}. Their sum is \eqref{shj:generalcontact}.
\end{proof}

The harmonic numbers in this theorem resemble the pointwise derivative coefficients in the canonical manuscript, but their role is different. Here they multiply delta derivatives introduced by periodic extension. The pointwise derivative identities remain correct; the correction is required only when transporting them into regularized periodic products.

\subsection{Every polygamma product in first Hurwitz derivatives}
For $r,k\ge0$ define
\[
\Pi_{r,k}(a)=\FP\int_0^1\psi^{(r)}(x)\psi^{(k)}(\{x+a\})\dd x.
\]
\begin{corollary}[Polygamma closure]\label{shj:polyclosure}
For $N=r+k\ge1$,
\begin{align}
\Pi_{r,k}(a)={}&(-1)^r\big[\gamma_1^{(N)}(a)-H_r\psi^{(N)}(a)\big]\notag\\
&+(-1)^k\big[\gamma_1^{(N)}(b)-H_k\psi^{(N)}(b)\big].
\label{shj:polyclosuregamma}
\end{align}
Equivalently,
\begin{align}
\Pi_{r,k}(a)=N!\bigg\{&(-1)^{k+1}\big[\zeta^{[1]}(N+1,a)
+(H_N-H_r)\zeta(N+1,a)\big]\notag\\
&+(-1)^{r+1}\big[\zeta^{[1]}(N+1,b)
+(H_N-H_k)\zeta(N+1,b)\big]\bigg\}.
\label{shj:polyclosurezeta}
\end{align}
The case $N=0$ is \eqref{shj:psipsi}.
\end{corollary}
\begin{proof}
Take $m=n=0$ in \eqref{shj:generalcontact}, use $\gamma_0^{(r)}=-\psi^{(r)}$, $h_{r,0}=H_r$, and differentiate \eqref{shj:J00}. This proves \eqref{shj:polyclosuregamma}. The pointwise identities
\begin{align*}
\psi^{(N)}(x)&=(-1)^{N+1}N!\zeta(N+1,x),\\
\gamma_1^{(N)}(x)&=(-1)^{N+1}N!
\big[\zeta^{[1]}(N+1,x)+H_N\zeta(N+1,x)\big]
\end{align*}
then give \eqref{shj:polyclosurezeta}. The second follows directly by expanding
$\partial_x^N\zeta(1+u,x)=(-1)^N(1+u)_N\zeta(N+1+u,x)$ at $u=0$; these pointwise ladders are also part of the repository's existing development \cite{shj:repo}.
\end{proof}

Two instructive special cases are
\begin{align}
\Pi_{0,1}(a)&=\zeta^{[1]}(2,a)+\zeta(2,a)-\zeta^{[1]}(2,b)
=J_{00}'(a)+\psi'(b),\label{shj:01}\\
\Pi_{1,1}(a)&=2\zeta^{[1]}(3,a)+2\zeta^{[1]}(3,b)
+\zeta(3,a)+\zeta(3,b).\label{shj:11}
\end{align}
The term $+\psi'(b)$ in \eqref{shj:01} is an explicit counterexample to the unqualified rule ``differentiate the finite-part identity under the integral.'' It is not a counterexample to ordinary differentiation under a genuinely dominated integral.

For reproducibility, let $I_{r,k}(\epsilon)$ be the sum of the two raw cutoff integrals, with arguments split as in \eqref{shj:Jcutoff}. Then
\begin{align}
\Pi_{0,1}(a)=\lim_{\epsilon\downarrow0}\left[
I_{0,1}(\epsilon)-\frac{\psi(b)}\epsilon
-\big(\psi'(a)-\psi'(b)\big)\log\epsilon\right],\label{shj:cut01}\\
\Pi_{1,1}(a)=\lim_{\epsilon\downarrow0}\left[
I_{1,1}(\epsilon)-\frac{\psi'(a)+\psi'(b)}\epsilon
+\big(\psi''(a)+\psi''(b)\big)\log\epsilon\right].\label{shj:cut11}
\end{align}
These counterterms can be derived directly from $\psi(x)=-1/x-\gamma+O(x)$ and $\psi'(x)=x^{-2}+O(1)$. They give numerical tests of the contact law that do not implement distributional differentiation.

\section{Ordinary integrals of all antiderivative orders}\label{shj:primitives}
The finite-part identities have a convergent counterpart. We first choose mean-zero periodic primitives, rather than implicitly choosing constants of integration at a singular endpoint.

\begin{definition}[Normalized primitive tower]\label{shj:Udef}
For integers $m\ge0$ and $r\ge1$, put
\begin{align}
\mathfrak h_j^{(r-1)}&=[u^j]\prod_{\nu=1}^{r-1}(1-u/\nu)^{-1},\notag\\
U_{m,r}(x)&=\frac{(-1)^{m+1}m!}{(r-1)!}
\sum_{j=0}^{m+1}\mathfrak h_{m+1-j}^{(r-1)}
\frac{\zeta^{[j]}(1-r,x)}{j!}.
\label{shj:Uformula}
\end{align}
The empty product convention gives $\mathfrak h_0^{(0)}=1$ and $\mathfrak h_j^{(0)}=0$ for $j>0$. Write $P_m=U_{m,1}$ and $Q_m=U_{m,2}$. Explicitly,
\begin{align}
P_m(x)&=\frac{(-1)^{m+1}}{m+1}\zeta^{[m+1]}(0,x),\label{shj:P}\\
Q_m(x)&=(-1)^{m+1}m!\sum_{j=0}^{m+1}\frac{\zeta^{[j]}(-1,x)}{j!}.
\label{shj:Q}
\end{align}
\end{definition}
These are normalized Hurwitz-jet primitives. The repository's primitive based at $1$ is $P_m(x)-P_m(1)$; it is not $P_m(x)$ itself. Likewise, zero-based negative polygamma conventions must be converted explicitly before using the following correlations.

\begin{lemma}[Primitive identities and regularity]\label{shj:Ulemma}
Each $U_{m,r}$ has mean zero on $(0,1)$. One has
\[
DP_m=G_m,\qquad DU_{m,r}=U_{m,r-1}\quad(r\ge2),\qquad D^rU_{m,r}=G_m
\]
as periodic distributions. The functions $P_m$ lie in $L^2(\T)$, and for $r\ge2$, $U_{m,r}$ is a periodic $C^{r-2}$ function.
\end{lemma}
\begin{proof}
Consider the spectral primitive
\begin{equation}\label{shj:primitivegenerator}
V_r(u,x)=\frac{(-1)^r\zeta(1+u-r,x)}{(1+u-r)_r}
=-\frac{\zeta(1+u-r,x)}{(r-1)!u}
\prod_{\nu=1}^{r-1}(1-u/\nu)^{-1}.
\end{equation}
Its regular coefficient $(-1)^m m![u^m]V_r$ is \eqref{shj:Uformula}. Pointwise differentiation gives $\partial_xV_r=V_{r-1}$ for $r\ge2$, and $\partial_xV_1=\zeta(1+u,x)$. Comparing regular coefficients proves the pointwise primitive relations.

For $r=1$, the shift identity gives
\[
P_m(x)=\frac{\log^{m+1}x}{m+1}+P_m(1+x).
\]
This proves square integrability. Integration by parts against a periodic test function, with a lower cutoff, yields the boundary term $\phi(0)\log^{m+1}\epsilon/(m+1)$ and hence $DP_m=G_m$ by \eqref{shj:Gdef}.
For $r\ge2$, the singular term in $\zeta^{[j]}(1-r,x)$ is $x^{r-1}(-\log x)^j$, which tends to zero at $x=0$. The endpoint values at $0+$ and $1$ consequently match. Iterating the pointwise derivative relations shows periodic $C^{r-2}$ regularity and no extra jump terms. This proves the distributional relations. Finally, $\int_0^1\zeta(s,x)\dd x=0$ for $\Re s<1$, and its spectral derivatives also integrate to zero. Apply this to each term in \eqref{shj:Uformula}.
\end{proof}

\begin{theorem}[All-order convergent lifting]\label{shj:alllift}
Suppose the closure formula for fixed $m,n$ is written as
\begin{equation}\label{shj:linearform}
J_{mn}(a)=\sum_j c_j\gamma_j(a)+\sum_j d_j\gamma_j(b)+c,
\end{equation}
where all sums are finite and the coefficients are independent of $a$.
For $r,k\ge1$, $N=r+k$, define the ordinary integral
\[
W_{m,n}^{r,k}(a)=\int_0^1U_{m,r}(x)U_{n,k}(\{x+a\})\dd x.
\]
Then
\begin{align}
W_{m,n}^{r,k}(a)={}&(-1)^r\sum_j c_jU_{j,N}(a)
+(-1)^k\sum_j d_jU_{j,N}(b)\notag\\
&+\frac{(-1)^r c}{N!}B_N(a),
\label{shj:allliftformula}
\end{align}
where $B_N$ is the ordinary Bernoulli polynomial. Thus every such integral reduces to finitely many Hurwitz-zeta jets at the single spectral integer $1-r-k$ and zeta-valued coefficients. The highest required spectral derivative is $m+n+2$.
\end{theorem}
\begin{proof}
The integral exists by the $L^2$ property and Cauchy--Schwarz. Its mean in $a$ is zero by Fubini and the mean-zero normalization of both factors. In the distributional sense, periodic integration by parts gives
\begin{equation}\label{shj:Wderivative}
\frac{d^N}{da^N}W_{m,n}^{r,k}(a)=(-1)^rJ_{mn}(a),\qquad 0<a<1.
\end{equation}
The right side is smooth there.

We need enough endpoint regularity to fix all integration constants. From \eqref{shj:fouriercoef} and \eqref{shj:Zlaurent}, for $q\ne0$,
\begin{align}
\widehat G_m(q)=(-1)^{m+1}m![u^{m+1}]
\exp\bigg[&\left(\kappa+\log|q|+\frac{\pi\ii}{2}\sgn q\right)u\notag\\
&+\sum_{j\ge2}\frac{\zeta(j)}j u^j\bigg].
\label{shj:Gfourier}
\end{align}
In particular $\widehat G_m(q)=O_m((1+\log|q|)^{m+1})$.
The primitive relations imply
$\widehat U_{m,r}(q)=\widehat G_m(q)/(2\pi\ii q)^r$.
The Fourier coefficients of $W$ are therefore
$O((1+\log|q|)^{m+n+2}/|q|^N)$.
After $N-2$ derivatives the series is still absolutely convergent. Hence $W$ is periodic $C^{N-2}$.

Let $F(a)$ be the proposed right side of \eqref{shj:allliftformula}. It also has mean zero and is periodic $C^{N-2}$, using Lemma~\ref{shj:Ulemma} and the endpoint identities for Bernoulli polynomials. Its $N$th derivative on $(0,1)$ is $(-1)^rJ_{mn}$: the reflected primitive contributes $(-1)^N$, and $B_N^{(N)}=N!$. Thus $W-F$ is a polynomial of degree at most $N-1$ on $(0,1)$. Matching its endpoint derivatives of orders $0,\ldots,N-2$ forces every nonconstant coefficient to vanish, starting with the leading one. Its zero mean then forces the remaining constant to vanish.

The closure uses indices at most $m+n+1$, and $U_{j,N}$ uses spectral derivatives at most $j+1$. This gives the final derivative-order assertion.
\end{proof}

\subsection{The first-primitive correlation}
For $K_{mn}(a)=W_{m,n}^{1,1}(a)$, the theorem simplifies to
\begin{equation}\label{shj:Klift}
K_{mn}(a)=-\sum_j c_jQ_j(a)-\sum_j d_jQ_j(b)-\frac c2B_2(a),
\qquad K_{mn}''(a)=-J_{mn}(a).
\end{equation}
Equivalently, ordinary differentiation of the master integral at $(s,t)=(0,0)$ gives
\begin{equation}\label{shj:KfromC}
K_{mn}(a)=(-1)^{m+n}m!n![s^{m+1}t^{n+1}]C(s,t;a).
\end{equation}
The signs and factorials follow from \eqref{shj:P}; no finite-part operation occurs in this integral.

For example, \eqref{shj:J11} gives the fully explicit convergent identity
\begin{align}
K_{11}(a)={}&-\tfrac12\big(Q_3(a)+Q_3(b)\big)
-2\zeta(2)\big(Q_1(a)+Q_1(b)\big)\notag\\
&-\zeta(3)\big(Q_0(a)+Q_0(b)\big)
+\tfrac12\big(\zeta(4)+\zeta(2)^2\big)B_2(a).
\label{shj:K11}
\end{align}
This evaluates a product of two second spectral derivatives at zero using derivatives at $-1$ of order at most four.

An example with unequal integration orders is
\begin{equation}\label{shj:W21}
\int_0^1Q_0(x)P_0(\{x+a\})\dd x
=U_{1,3}(a)-U_{1,3}(b)-\frac{\zeta(2)}3B_3(a),
\end{equation}
where
\[
U_{1,3}(x)=\frac78\zeta(-2,x)+\frac34\zeta^{[1]}(-2,x)
+\frac14\zeta^{[2]}(-2,x).
\]
Thus the same finite closure gives exact identities involving successive antiderivatives rather than merely a formal recurrence for integration constants.

\section{Explicit log-gamma and mixed identities}\label{shj:loggamma}
Let
\[
f(x)=\log\Gamma(x)-\frac\ell2,\qquad
\mathcal K(a)=\int_0^1 f(x)f(\{x+a\})\dd x.
\]
Lerch's identity gives $P_0=-f$, so $\mathcal K=K_{00}$. No regularization is used in this definition: logarithmic singularities are square integrable.

\begin{theorem}[Translated log-gamma correlation]\label{shj:gammathm}
For $0<a<1$,
\begin{align}
\mathcal K(a)={}&(1+\zeta(2))B_2(a)
-\zeta^{[1]}(-1,a)-\zeta^{[1]}(-1,b)\notag\\
&-\frac12\big[\zeta^{[2]}(-1,a)+\zeta^{[2]}(-1,b)\big].
\label{shj:gammahurwitz}
\end{align}
Equivalently, for $z=e^{2\pi\ii a}$,
\begin{equation}\label{shj:gammapoly}
\mathcal K(a)=\frac1{2\pi^2}\Re\left[
\Li_2^{[2]}(z)-2\kappa\Li_2^{[1]}(z)
+\left(\kappa^2+\frac{\pi^2}{4}\right)\Li_2(z)\right].
\end{equation}
For the uncentered product,
\begin{equation}\label{shj:gammauncentered}
\int_0^1\log\Gamma(x)\log\Gamma(\{x+a\})\dd x
=\frac{\ell^2}{4}+\mathcal K(a).
\end{equation}
Finally,
\begin{equation}\label{shj:gammasecond}
\mathcal K''(a)=\frac{\pi^2}{3}-\gamma_1(a)-\gamma_1(b).
\end{equation}
\end{theorem}
\begin{proof}
Apply \eqref{shj:Klift} to $J_{00}=\gamma_1(a)+\gamma_1(b)-2\zeta(2)$. Since
\[
Q_1(x)=\zeta(-1,x)+\zeta^{[1]}(-1,x)+\tfrac12\zeta^{[2]}(-1,x)
\]
and $\zeta(-1,x)=-B_2(x)/2$, the result is \eqref{shj:gammahurwitz}.

For an independent Fourier derivation of \eqref{shj:gammapoly}, differentiate \eqref{shj:fouriercoef} at $s=0$. The positive-frequency coefficient of $f$ is
\[
\widehat f(n)=\frac1{4n}-\frac{\ii(\kappa+\log n)}{2\pi n},\qquad n\ge1.
\]
The negative-frequency coefficient is its conjugate. These coefficients also give the classical Kummer Fourier expansion. The $L^2$ correlation therefore equals the absolutely convergent series
\begin{equation}\label{shj:gammafourier}
\mathcal K(a)=\frac1{2\pi^2}\sum_{n\ge1}
\frac{\cos(2\pi na)}{n^2}
\left[(\kappa+\log n)^2+\frac{\pi^2}{4}\right].
\end{equation}
Taking the first two spectral derivatives of the absolutely convergent polylogarithm series at order $2$ proves \eqref{shj:gammapoly}. The mean of $f$ is zero, so expanding the uncentered product proves \eqref{shj:gammauncentered}. Equation \eqref{shj:gammasecond} is the $m=n=0$ case of \eqref{shj:Klift}.
\end{proof}

\begin{corollary}[A mixed log-gamma--digamma identity]\label{shj:mixedgamma}
The same finite-part convention gives
\begin{align}
\FP\int_0^1\log\Gamma(x)\psi(\{x+a\})\dd x
={}&\frac12\big[\zeta^{[2]}(0,b)-\zeta^{[2]}(0,a)\big]\notag\\
&+\zeta(2)(2a-1).
\label{shj:mixedformula}
\end{align}
It is unchanged when $\log\Gamma(x)$ is replaced by $f(x)$.
\end{corollary}
\begin{proof}
Since $DP_0=G_0$ with no contact term at this first primitive stage,
$\mathcal K'(a)=\FP\int f(x)\psi(\{x+a\})\dd x$.
Differentiate \eqref{shj:gammahurwitz} using
\[
\partial_x\zeta^{[j]}(-1,x)
=\zeta^{[j]}(0,x)-j\zeta^{[j-1]}(0,x).
\]
The first spectral derivatives at zero cancel, and the terms from
$\zeta(0,x)=1/2-x$ simplify to $\zeta(2)(2a-1)$. This gives \eqref{shj:mixedformula}. The canonical periodic finite part of $\psi$ has mean zero, so subtracting $\ell/2$ from the log-gamma factor does not change the answer.
\end{proof}
For example, its left side with the centered function equals the entirely convergent expression
\begin{align*}
&\int_0^b f(x)\psi(x+a)\dd x\\
&\quad+\int_0^a\left[(f(b+t)-f(b))\psi(t)+f(b)\psi(1+t)\right]\dd t
-f(b)\log a.
\end{align*}
The recurrence $\psi(1+t)=\psi(t)+1/t$ removes the endpoint pole explicitly. This form also provides a direct numerical check independent of differentiating a finite-part integral.

\subsection{Half shift and the zero-shift limit}
At $a=1/2$, use $\Li_s(-1)=(2^{1-s}-1)\zeta(s)$. With $L_2=\log2$,
\begin{align}
\mathcal K(1/2)={}&-\frac{\zeta^{[2]}(2)}{4\pi^2}
+\frac{\kappa-L_2}{2\pi^2}\zeta^{[1]}(2)\notag\\
&+\frac{L_2^2+2\kappa L_2-\kappa^2-\pi^2/4}{24}.
\label{shj:gammahalf}
\end{align}
This formula concerns a \emph{translated} correlation, not the reflected product $f(x)f(1-x)$.

The absolutely convergent series \eqref{shj:gammafourier} is continuous at $a=0$. Its value there recovers the classical second log-gamma moment \cite{shj:EM}:
\begin{align}
\int_0^1\log^2\Gamma(x)\dd x={}&
\frac{\ell^2}{4}+\frac{\zeta^{[2]}(2)}{2\pi^2}
-\frac{\kappa\zeta^{[1]}(2)}{\pi^2}
+\frac{\kappa^2}{12}+\frac{\pi^2}{48}.
\label{shj:classicalmoment}
\end{align}
Equivalently, taking the continuous limit in \eqref{shj:gammahurwitz} gives
\[
\mathcal K(0)=\frac{1+\zeta(2)}6-2\zeta^{[1]}(-1)-\zeta^{[2]}(-1).
\]
These are useful consistency checks, not priority claims. In contrast,
$J_{00}(a)=\log a/a+O(1)$ as $a\downarrow0$, by \eqref{shj:localgamma}. The finite-part product has no continuous zero-shift specialization. The ordinary correlation remains continuous but its second derivative develops the expected singularity.

\section{Rational shifts and exact arithmetic coordinates}\label{shj:rationals}
For integers $q\ge2$ and $1\le p<q$, residue-class decomposition gives the classical finite Fourier identity
\begin{equation}\label{shj:DFT}
\Li_s(e^{2\pi\ii p/q})
=q^{-s}\sum_{r=1}^q e^{2\pi\ii pr/q}\zeta(s,r/q).
\end{equation}
It is proved by absolutely convergent series for $\Re s>1$ and then by analytic continuation. Spectral differentiation yields
\begin{align}
\Li_s^{[j]}(e^{2\pi\ii p/q})
=q^{-s}\sum_{r=1}^q e^{2\pi\ii pr/q}
\sum_{h=0}^j\binom jh(-\log q)^{j-h}\zeta^{[h]}(s,r/q).
\label{shj:DFTjets}
\end{align}
At $s=1$ the poles of the finite sum cancel because the roots sum to zero. One must combine the pole terms before evaluating; separately substituting $\zeta(1,r/q)$ is invalid.

At a rational shift, Theorem~\ref{shj:closurethm} reduces all Stieltjes correlations to rational-argument Stieltjes constants and ordinary zeta values. Theorem~\ref{shj:alllift} places every primitive correlation at a negative integer spectral argument. Equation \eqref{shj:DFTjets} instead evaluates the polylogarithmic side using positive-order Hurwitz jets. The two routes give independent choices of arithmetic coordinates for numerical replay and exact symbolic comparison.

For a Dirichlet character modulo $q$, including an imprimitive one with its actual modulus, the identity
\[
\sum_{r=1}^q\chi(r)\zeta(s,r/q)=q^sL(s,\chi)
\]
and its spectral derivatives transform the character components into Dirichlet-$L$ jets. This gives a reduction framework, not a claim of numerical independence or minimality of the resulting coordinates.

\begin{corollary}[Small-denominator digamma correlations]\label{shj:rationalpsi}
Write $\gamma_1=\gamma_1(1)$. Then
\begin{align}
\Pi_{0,0}(1/2)&=2\gamma_1-4\gamma\log2-2\log^22-\frac{\pi^2}{3},\label{shj:rathalf}\\
\Pi_{0,0}(1/3)&=2\gamma_1-3\gamma\log3-\frac32\log^23-\frac{\pi^2}{3},\label{shj:ratthird}\\
\Pi_{0,0}(1/4)&=2\gamma_1-6\gamma\log2-7\log^22-\frac{\pi^2}{3}.
\label{shj:ratquarter}
\end{align}
\end{corollary}
\begin{proof}
Expand at $s=1$ the exact identities
\begin{align*}
\zeta(s,1/2)&=(2^s-1)\zeta(s),\\
\zeta(s,1/3)+\zeta(s,2/3)&=(3^s-1)\zeta(s),\\
\zeta(s,1/4)+\zeta(s,3/4)&=(4^s-2^s)\zeta(s).
\end{align*}
Compare the coefficient of $s-1$ with \eqref{shj:laurent} and insert the resulting sums into \eqref{shj:psipsi}. In the first line the correlation contains two copies of $\gamma_1(1/2)$, which accounts for its different coefficients.
\end{proof}
Higher indices require no new search: expand these finite exponential factors to the needed order and use \eqref{shj:closureformula}. At larger denominators, retaining a nontrivial character block is legitimate; failing to reduce that block does not establish transcendence or independence.

\section{Verification and reproducible computation}\label{shj:verification}
The analytic proofs above are the justification of the identities. The computational artifacts serve three different purposes: exact coefficient replay, independent numerical comparison, and reproducible document production. None is described as a proof-assistant formalization or as interval certification.

\subsection{Exact replay}
The script \src{code/identity_engine.py} works over rational polynomials in formal symbols $\zeta(2),\zeta(3),\ldots$. It emits all $28$ ordered pairs $(m,n)$ with $m+n\le6$ and their first-primitive lifts. It checks exact divisibility in \eqref{shj:AandT}, the symmetry $J_{mn}(a)=J_{nm}(1-a)$, linearity in the Stieltjes symbols, and weight conservation. It also checks the harmonic contact coefficients for derivative orders through eight and $35$ primitive recurrences, for $2\le r\le6$, $0\le m\le6$.

The implementation truncates \emph{total} degree rather than rectangular degree, because the quotient kernel lowers degree by two. Equation \eqref{shj:coefficientalgorithm} avoids expanding a large product containing every Stieltjes symbol. The JSON output records exact rational coefficients and readable expressions, so a different computer algebra system can replay the identities without a numerical tolerance.

\subsection{Numerical diagnostics}
The numerical script uses Python with mpmath at $50$ decimal digits. Its checked representations include defining ordinary integrals, the two-variable beta formula, a separately implemented polylogarithmic generator at $a=1/2$, direct endpoint-cutoff integrals, and Barnes-$G$ evaluations of a higher primitive. The completed run contains the following $64$ checks:
\begin{center}
\begin{tabular}{@{}p{0.76\linewidth}r@{}}
\toprule
Diagnostic family & Count\\
\midrule
Euler--Maclaurin Stieltjes evaluator versus mpmath & 12\\
Master correlation: Fourier, beta, and defining integrals & 4\\
Independent polylogarithmic jets through total index six & 28\\
Log-gamma, its polylogarithmic form, and mixed digamma integrals & 9\\
Higher primitive correlations, orders $(2,1)$ and $(2,2)$ & 2\\
Rational digamma specializations & 3\\
Direct Stieltjes finite parts at cutoff $10^{-18}$ & 4\\
Direct polygamma contact-term tests at cutoff $10^{-18}$ & 2\\
\bottomrule
\end{tabular}
\end{center}
All pass their stated tolerances. In particular, the four finite-cutoff Stieltjes tests at $a=0.3$ have residuals approximately
\[
1.24\times10^{-17},\quad 5.22\times10^{-16},\quad
1.15\times10^{-15},\quad 2.21\times10^{-14}
\]
for $J_{00},J_{10},J_{11},J_{20}$, respectively. Their common acceptance threshold is $10^{-13}$; the remaining discrepancy includes the intentional finite-cutoff error. These residuals are not asserted to be enclosures of an infinite-limit computation.

The Stieltjes integral evaluator uses an Euler--Maclaurin expansion with $96$ initial summands and $40$ correction terms. Its local checks compare against \src{mpmath.stieltjes}. For endpoint quadrature the substitution $x=c e^{-t}$ turns powers of $\log x$ into ordinary polynomial growth over a finite interval. The mixed log-gamma test removes the endpoint pole analytically before quadrature. The higher-primitive check evaluates
\[
Q_0(x)=\frac{B_2(x)}2-\zeta^{[1]}(-1)
+\log G(x+1)-x\log\Gamma(x)
\]
using the classical Barnes-$G$/Hurwitz relation, whereas its proposed answer uses Hurwitz jets at $-2$ or $-3$.

A numerical differentiation warning arose during development: naive tiny-step differentiation of the generic complex polylogarithm near integer order, and of a generic Hurwitz call at order zero, lost digits in the tested implementation. The retained script instead uses the explicit Hurwitz derivative argument and exact rational Fourier decomposition for the polylogarithm jets at order two. This is a diagnostic observation about the tested calls, not a mathematical counterexample or a general claim that the library's special functions are inaccurate.

The authoritative machine-readable record is \src{data/numerical_validation.json}, including software versions, every residual and tolerance, and the execution duration. The complete command is
\begin{verbatim}
python code/identity_engine.py --max-total 6
python code/verify_numerically.py --dps 50
python code/build.py
\end{verbatim}
The optional \src{--quick} numerical mode omits the four expensive Stieltjes cutoff quadratures; the delivered full record does not use that option.

\section{Repository audit and integration boundaries}\label{shj:audit}
The inspection focused on the canonical integration and differentiation chapters, the reflected-moment continuation, the discovery/status chapter, the manuscript README, the incoming inventory and intake protocol, and the observed branch head. It was not a line-by-line audit of the entire book or every incoming archive. The provenance file records this limitation rather than treating an archive listing as a scientific review of its contents.

\subsection{A concrete wording correction}
The exploratory note labeled \src{integral:neg:psim2} in the integration chapter says that basis atoms must be $\Q$-independent. Taken as a universal prerequisite, this is too strong: unknown independence is neither required to search for a relation nor required to prove a proposed relation by an exact identity. Known dependencies can instead be removed, quotiented out, or explicitly separated from target-bearing relations. A precise replacement is:
\begin{quote}
Remove or quotient known dependencies in the search basket, and distinguish relations involving the target from pre-existing relations among the basket atoms. Unknown numerical independence is not a prerequisite for an exact proof.
\end{quote}
This is a methodological wording correction. It does not invalidate the surrounding exact gamma-integral formulas, and it is consistent with the manuscript's later, more careful discussion of integer-relation searches.

\subsection{A boundary warning, not a correction to a true pointwise theorem}
The manuscript's pointwise identities
\[
\partial_x^r\zeta(s,x)=(-1)^r(s)_r\zeta(s+r,x)
\]
and its Stieltjes derivative tower are correct in their stated pointwise setting. They must not be transported into canonical periodic finite parts by silently assuming that differentiation commutes with extension. The correct replacement in that setting is Theorem~\ref{shj:contactthm}. For example, omitting the $+\psi'(1-a)$ term in \eqref{shj:01} gives a false identity. This report proposes adding an extension warning; it does not claim that the inspected manuscript already asserts that false periodic identity.

Other integration safeguards are equally concrete. Keep the finite-part coordinates fixed. Combine the beta terms before evaluating at a removable singularity. Do not substitute the collision point $a=0$ into separated-singularity formulas. Convert base-point-normalized antiderivatives to the mean-zero primitives before applying the lifting theorem. Treat exact symbolic closure as a representation theorem, not as period independence.

\subsection{Placement and acceptance}
A proposed report destination is
\src{Analysis/Polylogarithms/docs/reports/shifted-hurwitz-jet-correlations/}.
The primitive identities naturally follow the existing integration chapter; the contact theorem naturally follows the pointwise derivative tower. Labels and bibliography keys in this article use the prefix \src{shj:} to facilitate extraction. The supplied source is standalone; the integration note lists the additional macros needed when importing its body into the book.

No repository files were modified during this delivery. The accompanying audit and proposed replacement note are review artifacts, not applied patches. In accordance with the incoming procedure \cite{shj:incoming}, preservation of this archive and successful code replay should remain separate from mathematical acceptance into the canonical manuscript.

\section{Further research questions}\label{shj:further}
\subsection{Three shifted factors and genuine higher depth}
For three distinct shifts, define the analogous finite-part product of three Stieltjes functions by local subtraction at each singular point. Fourier integration imposes one relation on three frequencies, leaving a double sum rather than the one-frequency sum underlying Theorem~\ref{shj:closurethm}. Determine a precise closed coordinate class: are colored Tornheim or multiple Lerch jets sufficient at every Laurent index? The concrete first target is a fully regularized formula for three digamma factors with pairwise distinct rational shifts. Linear closure in single Stieltjes functions should not be assumed.

\subsection{Collision and renormalized same-point products}
The functions $J_{mn}(a)$ become singular as $a\to0$. Their Fourier-defined convolutions nevertheless exist as distributions on the whole circle. Determine the complete local expansion, including delta derivatives, and compare its finite constant with a separately defined same-point Hadamard product. A useful target is an explicit conversion law between these two regularizations at every $(m,n)$, rather than an unjustified substitution into the separated formula.

\subsection{Dilation, traces, and subtraction-scale dependence}
Integer dilation $x\mapsto\{qx\}$ transports an endpoint coordinate to $qx$ and changes logarithmic counterterms by powers of $\log q$. Determine the exact commutators between dilation, canonical extension, and the derivative operation. This should produce a compatible distributional version of the manuscript's pointwise distribution and conductor laws. The question is about explicit identities and local correction coefficients, not about numerical rank inferred from a truncated matrix.

\subsection{A route toward fixed-weight polylogarithm conjectures}
The present formulas involve spectral derivatives of classical polylogarithms and bilinear Hurwitz correlations. They do not automatically reduce the multiple-polylogarithm periods occurring in $S_6$ and $S_8$. A concrete next step is to seek a proved integral transformation expressing a candidate's obstruction in a shifted two- or three-factor correlation algebra, with all endpoint terms retained. Only after such a bridge is established could the closure or a higher-depth extension become an exact certificate for those candidates.

\subsection{Mixed orders and further explicit gamma identities}
Theorem~\ref{shj:alllift} gives a uniform answer for integer integration orders. Develop convenient Barnes multiple-gamma coordinates for $U_{m,r}$ at small $m$ and large $r$, and compare them with the zero-based negative-polygamma convention used elsewhere in the repository. Explicit rational-shift tables in these coordinates would turn the general theorem into useful gamma and multiple-gamma product-integral identities. Analytic interpolation in the integration order is a separate question requiring a specified fractional periodic primitive.

\subsection{Formal certificates with analytic hypotheses attached}
The finite algebra in \eqref{shj:coefficientalgorithm} and \eqref{shj:hdef} is suitable for formal verification. A useful first formalization would certify the coefficient theorem conditional on a precisely stated meromorphic correlation identity, followed by the analytic Fourier and finite-part lemmas. A certificate should record the shift domain, normalization, derivative convention, and every contact term; a zero residual or a formally correct polynomial expansion without those hypotheses is not a proof of the integral identity.

\section*{Conclusion}
\addcontentsline{toc}{section}{Conclusion}
A translated Hurwitz product is an exact bridge between polylogarithm order derivatives and Stieltjes functions. Near the double pole, one divisible gamma quotient forces an all-order linear closure. Periodic differentiation adds computable harmonic contact terms, while normalized integration converts the same closure into ordinary convergent identities at negative spectral integers. The resulting framework supplies proved families of exact identities and a clear boundary between classical pointwise calculus, finite-part extension, and unresolved arithmetic reductions.
